\section{本讲小结}

无限网络首先是有限截断和边界条件的极限。自相似法应检查切口端点、首单元是否包含回流导轨，以及缩放是否同时作用于电阻。代数根须符合正性，并由截断序列或空间衰减证明其物理意义。

节点法给出的差分方程揭示电流随距离的衰减，源节点守恒确定系数，远处边界筛选特征根。有限网络一般需要保留两根，不能提前套用半无限衰减条件。虚拟电流法通过叠加抵消非目标节点的外部注入，在晶格中尤其方便。

相邻电阻的均分结论依赖具体网络的平移与局部旋转对称，不能只看节点度数。自相似三端网络需保留端口对应，必要时使用星形等效。最后通过已知特例、增删边单调性、极限和电流总和检查整个建模过程。

